\subsection{Pure AC: Positive Images and Finite Complete Families}
\label{sec:p12-equational-ac}

The source distinguishes AC from ACU and refers to the AC algorithms
of Stickel and others; its semantic \S5 develops the monoidal examples
ACU, ACUI and AG, not an implicit unit for AC. The variable-only AC
subset construction below is explicitly present in the separately
identified 1975 paper. Its extension here to simultaneous equations and
fixed free constants has full algebraic proofs and a finite basis algorithm.
It does not identify that 1975 manifestation with the 1981 journal paper.
\cite[\S3.4 and \S5.1, pp.~484--486,504--508]{BaaderSnyder2001}
\cite[pp.~72--75]{Stickel1975AC}

\begin{hbdefinition}{def:p12-ac-input}{The unit-free term domain and input matrices}
Let \(AC=\{f(f(x,y),z)\approx f(x,f(y,z)),f(x,y)\approx f(y,x)\}\),
with binary f and no designated unit or unit axiom. A flat AC problem is a finite set of
equations whose terms use f, variables, and free constants only; nested
f nodes are flattened. A single variable or constant is a legal term.
An empty product is not. In this construction no further positive-arity
symbol is in the term signature. In examples, multiargument f notation
abbreviates a binary fold, not a new arity; a one-atom product is that
atom itself. All substitution images are finite nonempty terms, and
constants are fixed. Compare answers only on the original
\(X=\operatorname{Var}(\Gamma)\), using the already declared AC-instance
preorder. Enumerate \(X=(x_1,\ldots,x_n)\), and the input constant
set C. Let \(M_L,M_R\) be the variable-multiplicity matrices and
\(u_{L,c},u_{R,c}\) each constant's multiplicity columns. Put
\(D=M_L-M_R\) and \(d_c=u_{R,c}-u_{L,c}\).
The homogeneous columns solve \(Dv=0,v\in\mathbb N^n\);
the c-columns solve \(Dw=d_c,w\in\mathbb N^n\).
Zero is an internal multiplicity vector, never an AC term image.
\end{hbdefinition}

\begin{hblemma}{thm:p12-ac-multisets}{proved=annex:p12.ac.multisets}{The exact unit-free coordinate model and cancellation}
For terms over f, variables and free constants, AC equality is equality
of their nonempty generator multisets. Substitution composition acts on
counts by nonnegative integer matrix multiplication. Common flattened
arguments can be cancelled without losing solutions. An empty residual
argument list after cancellation is only an internal balance constraint;
it neither introduces a unit nor changes the nonempty-image condition.
\cite[pp.~72,74]{Stickel1975AC}
\end{hblemma}

\begin{hbmethod}{def:p12-ac-single-bound}{The source's finite single-equation basis search}
After cancelling common variable counts, a single homogeneous equation
has positive coefficients on its remaining two sides:
\(\sum_{i=1}^{m}a_ix_i=\sum_{j=1}^{n}b_jy_j=k\).
If both sides are present, let
\(L=\max_{i,j}\operatorname{lcm}(a_i,b_j)\) and
\(K=\max(m,n)L\). Enumerate all nonzero nonnegative solutions
with balance k at most K; each coordinate is bounded by K divided
by its coefficient. Retain only the componentwise minimal solutions.
Variables cancelled completely are unconstrained, not assigned zero.
A remaining one-sided positive equation has only the zero solution.
\cite[Lemma and algorithm, pp.~73--75]{Stickel1975AC}
\end{hbmethod}

\begin{hblemma}{thm:p12-ac-single-bound}{proved=annex:p12.ac.single-bound}{The actual lcm bound proves finite generation}
Every nonnegative solution of the positive-coefficient equation in
\cref{def:p12-ac-single-bound} is an additive combination of solutions
with k at most K. Every minimal nonzero solution is therefore within
that finite search. The retained solutions generate all solutions.
\end{hblemma}

\begin{hbmethod}{def:p12-ac-system-basis}{A finite simultaneous-equation basis algorithm}
For any integer matrix A with n columns, enumerate every nonempty column
support S. Compute \(\ker A_S\) by rational Gaussian elimination.
If its dimension is one and a spanning vector has all coordinates
of one strict sign, choose its positive primitive integer vector,
extended by zeros off S. Collect the distinct vectors \(r_1,\ldots,r_t\).
Put \(B_i=\sum_j(r_j)_i\), taking B zero if no vector exists.
Enumerate the finite integer box \(0\leq v_i\leq B_i\), retaining
its nonzero solutions of \(Av=0\). Remove any vector having a
distinct nonzero retained solution componentwise below it; call the
remaining set \(H(A)\). All operations use exact rational/integer
arithmetic. This is an editorial system-level implementation, not a
claim that this ray/box procedure is Stickel's printed single-equation
search or the source's cited incremental algorithms.
\end{hbmethod}

\begin{hbtheorem}{thm:p12-ac-system-basis}{proved=annex:p12.ac.system-basis}{The finite search returns exactly the minimal homogeneous generators}
The preceding algorithm terminates and returns precisely the minimal
nonzero nonnegative integer solutions of \(Av=0\). They generate
all nonnegative integer solutions by addition. The ray set alone may
not generate the integer solutions; the bounded integer remainder
search is essential. No unbounded enumeration with an unstated stopping
criterion is used.
\end{hbtheorem}

\begin{hbmethod}{def:p12-ac-particulars}{Fixed constants through an augmented homogeneous system}
For \(Dw=d_c\), apply the basis algorithm to
\(\widetilde D_c=[D\mid -d_c]\). Select the basis vectors whose
last coordinate is one and project their first n coordinates; this
finite set is \(P_c\). Compute \(H=H(D)=\{h_1,\ldots,h_k\}\)
once. A missing \(P_c\) signals that no c-column is possible.
\end{hbmethod}

\begin{hblemma}{thm:p12-ac-particulars}{proved=annex:p12.ac.particulars}{All and only minimal affine columns are found}
The projected set \(P_c\) is exactly the componentwise minimal
nonnegative solutions of \(Dw=d_c\). Every solution is
\(p+v\), with \(p\in P_c\) and v a homogeneous solution.
This includes \(d_c=0\), where \(P_c=\{0\}\); constants
cancelled from the equations can still occur in later parameter images.
\end{hblemma}

\begin{hbmethod}{def:p12-ac-extraction}{Positive-support branch construction and canonical extraction}
Choose one \(p_c\in P_c\) for each input constant. Enumerate
every subset \(I\subseteq\{1,\ldots,k\}\). Retain it only if
\[
             \sum_{c\in C}(p_c)_i+\sum_{j\in I}(h_j)_i>0
                    \qquad(1\leq i\leq n).
\]
Use a distinct fresh variable \(z_j\) for each selected j. Its
image in \(x_i\sigma_{p,I}\) has \((h_j)_i\) occurrences,
and c has \((p_c)_i\) occurrences. Assemble that nonempty multiset
as one atom when its size is one and as a fixed right-associated
binary f term otherwise. All variables outside X are fixed.
The empty I is permitted when the constants alone make every image
nonempty. A zero row is rejected, never represented by an invented
unit or \(f()\). An input constant named e is an ordinary atom,
not a zero-multiplicity image.
Canonical integer vectors, fixed original-variable order, sorted
constants and fixed fresh-variable choices deduplicate literal outputs.
For a supplied target unifier, its factor can also be computed:
choose the first \(p_c\leq w_c\) in each finite \(P_c\), then
repeatedly subtract the first \(h_j\leq v\) from each homogeneous
residual column v until it is zero. Record the subtraction counts as
parameter-image multiplicities. Each subtraction strictly lowers the
coordinate sum, so this factor construction terminates as well.
\end{hbmethod}

\begin{hbtheorem}{thm:p12-ac-coverage}{proved=annex:p12.ac.coverage}{The complete finite pure-AC algorithm}
For a flat finite AC problem, the basis, augmented-column and positive-
support procedures terminate and return a finite complete set of AC
unifiers. Every arbitrary AC unifier is covered by a returned answer
whose factor assigns a nonempty term to each selected fresh variable.
Zero coefficients decide which generators are omitted, not which
parameters are assigned empty terms. Empty input returns identity.
For n=0 compare the ground input multisets directly, returning identity
if they agree and the empty answer set otherwise.
\end{hbtheorem}

\begin{hbtheorem}{thm:p12-ac-minimality}{proved=annex:p12.ac.minimality}{Minimality and deduplication are established, not assumed}
With the actual minimal H and minimal \(P_c\) computed above,
distinct admissible pairs \((p,I)\) give incomparable answers
on X, up to the fixed canonical fresh naming. Thus the deduplicated
output is a minimal complete set, not merely complete. A nonminimal
generating set would not justify this conclusion. No uniqueness of
factor substitutions is asserted.
\end{hbtheorem}

\begin{hbexample}{def:p12-ac-five-answers}{The source's ACU equation has five unit-free answers}
For \(f(x,f(x,y))=^{?}_{AC}f(z,f(z,z))\), H consists of
\(h_1=(1,1,1),h_2=(0,3,1),h_3=(3,0,2)\).
Exactly five subsets cover every variable row:
\(\{1\},\{2,3\},\{1,2\},\{1,3\},\{1,2,3\}\).
For example \(\{1\}\) yields x=y=z=\(z_1\), whereas
\(\{2,3\}\) yields \(x=f(z_3,f(z_3,z_3))\),
\(y=f(z_2,f(z_2,z_2))\), \(z=f(z_2,f(z_3,z_3))\).
Singleton 2 is rejected for its empty x-image; singleton 3 for its
empty y-image. The five answers are minimal complete by the theorem.
ACU instead allows switching off parameters by e and combines them
into one answer. ACUI changes counts to supports, and AG allows signed
counts; neither operation belongs to this AC algorithm.
The two positive cone rays are h2 and h3; h1 is an indispensable
integer remainder, not an integer sum of those rays.
\cite[pp.~505--506]{BaaderSnyder2001}
\end{hbexample}

\begin{hbexample}{def:p12-ac-stickel-69}{The actual 1975 seven-generator and constant example}
The source variable equation is
\(2x_1+x_2+x_3=2y_1+y_2\). In source order its seven minimal
vectors are
\[
\begin{gathered}
(0,0,2,1,0),\ (0,1,0,0,1),\ (0,0,1,0,1),\ (0,1,1,1,0),\\
(0,2,0,1,0),\ (1,0,0,0,2),\ (1,0,0,1,0).
\end{gathered}
\]
Exactly 69 of its 128 subsets have positive sums in all five rows.
For \(f(x,x,y,a)=^{?}_{AC}f(b,b,z)\), constant a is the
third row and b the fourth. A constant row must contain exactly
one parameter occurrence, not an f-product of several occurrences.
This leaves six subsets; any using vector 4 forces the same parameter
to equal distinct a and b, leaving four answers. The direct fixed-
constant procedure computes the same minimal family, with fresh p,q:
\[
\begin{aligned}
 &(x,y,z)=(b,p,f(a,p)),\\
 &(x,y,z)=(f(b,q),p,f(a,f(p,f(q,q)))),\\
 &(x,y,z)=(q,f(b,b),f(a,f(q,q))),\\
 &(x,y,z)=(q,f(b,f(b,p)),f(a,f(p,f(q,q)))).
\end{aligned}
\]
Here \(H(D)=\{(0,1,1),(1,0,2)\}\),
\(P_a=\{(0,0,1)\}\), and
\(P_b=\{(1,0,0),(0,2,0)\}\). These yield the four positive
branches without treating a or b as substitutable variables.
\cite[pp.~73--75]{Stickel1975AC}
\end{hbexample}

\begin{hbexample}{def:p12-ac-degeneracies}{Cancellation, constants and impossible positive images}
For \(x=^{?}_{AC}f(x,x)\), the homogeneous kernel is zero;
its only numerical column would make the image empty, so the answer
set is empty. ACU has \(x\mapsto e\); idempotent ACUI has identity.
For \(f(x,a)=^{?}_{AC}f(x,b)\) with a,b distinct, cancellation
leaves incompatible fixed constant columns, so there is no answer.
For \(f(x,y)=^{?}_{AC}f(a,b)\), the two answers are
\((x,y)=(a,b)\) and \((b,a)\). The other numerical allocations
leave one image empty and are rejected. ACU additionally permits
\((f(a,b),e)\) and \((e,f(a,b))\).
An identically zero variable column is unconstrained and has a unit
coordinate generator; it must still receive a nonempty variable image.
For \(f(x,x)=^{?}_{AC}a\), the augmented kernel has generator
\((1,2)\), but no last-coordinate-one generator; there is no
integer particular column, even though a rational half-column exists.
Duplicating an equation does not change the kernel or answer classes.
All these cases are handled by the stated finite procedure.
\end{hbexample}

\begin{proofboundary}{def:p12-ac-general-boundary}{Flat construction and general-input source boundaries}
This procedure covers every flat f/variable/constant system. It does
not silently turn a free-function-rooted input argument into a constant.
Stickel 1975 generalizes such arguments to fresh variables and then
reconciles their values by recursive unification, imposing single-atom
constraints where required. Its general-term proof invokes recursive
calls of lesser complexity. The paper expressly states termination,
soundness and completeness for that general-term algorithm; the flat
proof is not presented as establishing that stronger source statement.
The historical recursive algorithm's unrestricted nested-input
termination argument and the distinct 1981 journal presentation are
not supplied by the flat matrix proof here. An independent full
route for arbitrary finite nested/multiple AC symbols and free
symbols/constants is supplied by the AC-LCR bridge and actual ordered
combination construction in
\cref{thm:p12-comb-ac-lcr,sec:p12-equational-combination}, with its
own complete proofs and terminating component procedures. That route
does not discharge these distinct historical-source presentations.
The source's \S5 similarly states that its semantic approach does
not directly solve general free-function inputs; the later combination
treatment supplies a justified bridge, not an unmentioned extension
of that semantic method.
\cite[pp.~73--76]{Stickel1975AC}
\cite[pp.~504,512--513]{BaaderSnyder2001}
\end{proofboundary}
